% ── LNCS document class (download llncs.cls from https://www.springer.com/gp/authors-editors/book-authors-editors/manuscript-preparation/5636) ──
\documentclass{llncs}

% ── Encoding & language ────────────────────────────────────────────────────
\usepackage[utf8]{inputenc}
\usepackage[T1]{fontenc}
\usepackage[spanish]{babel}

% ── Mathematics ────────────────────────────────────────────────────────────
\usepackage{amsmath, amssymb}

% ── Code listings ──────────────────────────────────────────────────────────
\usepackage{listings}
\usepackage{xcolor}

\definecolor{codebg}{RGB}{248,248,248}
\definecolor{codecomment}{RGB}{106,153,85}
\definecolor{codekw}{RGB}{86,156,214}
\definecolor{codestr}{RGB}{206,145,120}

\lstdefinestyle{python}{
  language=Python,
  backgroundcolor=\color{codebg},
  basicstyle=\ttfamily\footnotesize,
  keywordstyle=\color{codekw}\bfseries,
  commentstyle=\color{codecomment}\itshape,
  stringstyle=\color{codestr},
  breaklines=true,
  frame=single,
  rulecolor=\color{gray!40},
  numbers=left,
  numberstyle=\tiny\color{gray},
  numbersep=8pt,
  showstringspaces=false,
  tabsize=4,
}

% ── Figures & tables ───────────────────────────────────────────────────────
\usepackage{graphicx}
\usepackage{booktabs}
\usepackage{array}

% ── Diagrams ───────────────────────────────────────────────────────────────
\usepackage{tikz}
\usetikzlibrary{shapes.geometric, arrows.meta, positioning, fit, backgrounds}

% ── Hyperlinks ─────────────────────────────────────────────────────────────
\usepackage[hidelinks, breaklinks]{hyperref}
\usepackage{url}

% ─────────────────────────────────────────────────────────────────────────
\begin{document}

% ── Title block ────────────────────────────────────────────────────────────
\title{Sistema de Recuperación de Información sobre Noticias Tecnológicas}

\author{Alberto}

\institute{Grado en Ingeniería Informática\\
\email{}}

\maketitle

% ── Abstract ───────────────────────────────────────────────────────────────
\begin{abstract}
Este trabajo presenta el diseño e implementación de un Sistema de Recuperación
de Información (SRI) sobre un corpus de noticias tecnológicas en español sobre
dispositivos móviles y ordenadores personales (PC). El sistema combina dos
modelos complementarios: BM25 para precisión léxica y LSI (\textit{Latent
Semantic Indexing}) para similitud semántica. Los documentos se obtienen
mediante arañas web (\textit{spiders}) construidas con Scrapy a partir de
prensa especializada. El preprocesado incluye tokenización, eliminación de
\textit{stopwords} y stemming morfológico. El procesador de consultas
incorpora expansión con sinónimos, inferencia de categoría y reordenación MMR.
Los índices se serializan a disco para permitir consultas sin reentrenamiento.
\end{abstract}

\keywords{recuperación de información \and BM25 \and LSI \and Scrapy \and
índice invertido \and procesamiento de lenguaje natural}

% ═══════════════════════════════════════════════════════════════════════════
\section{Definición del Dominio Seleccionado}
% ═══════════════════════════════════════════════════════════════════════════

\subsection{Descripción general}

El dominio seleccionado para este sistema de recuperación de información (SRI)
es el de \textbf{noticias tecnológicas}, concretamente artículos de prensa
especializada sobre dispositivos móviles y ordenadores personales (PC).

La colección de documentos se extrae de forma automática mediante arañas web
(\textit{spiders}) de Scrapy a partir de la siguiente fuente:

\begin{table}[h]
  \centering
  \caption{Fuente de datos del corpus tecnológico.}
  \begin{tabular}{lll}
    \toprule
    \textbf{Fuente} & \textbf{Idioma} & \textbf{Categoría} \\
    \midrule
    Xataka (\texttt{xataka.com}) & Español & Móvil y PC \\
    \bottomrule
  \end{tabular}
\end{table}

\subsection{Características del corpus}

El corpus presenta las siguientes particularidades que condicionan el diseño
del sistema:

\begin{itemize}
  \item \textbf{Riqueza sinonímica}: los textos tecnológicos emplean
    numerosos sinónimos y acrónimos: procesador/chip/CPU,
    pantalla/display/screen, cámara/fotografía/camera.
  \item \textbf{Polisemia}: términos como ``galaxy'' (Samsung) o ``apple''
    (empresa vs.\ fruta) exigen disambiguación por contexto.
  \item \textbf{Contenido mixto}: los artículos incluyen análisis, noticias,
    comparativas y tutoriales.
  \item \textbf{Metadatos ricos}: cada documento dispone de título, contenido
    completo, autor, fecha de publicación, etiquetas temáticas, marca y
    sistema operativo.
\end{itemize}

\subsection{Esquema canónico de documento}

Tras la extracción y normalización por el módulo \texttt{DocumentStore}, cada
documento sigue el esquema:

\begin{table}[h]
  \centering
  \caption{Esquema canónico de documento en \texttt{DocumentStore}.}
  \begin{tabular}{lll}
    \toprule
    \textbf{Campo} & \textbf{Tipo} & \textbf{Descripción} \\
    \midrule
    \texttt{id}           & str       & URL del artículo (identificador único) \\
    \texttt{title}        & str       & Titular del artículo \\
    \texttt{content}      & str       & Cuerpo completo del texto \\
    \texttt{author}       & str       & Autor \\
    \texttt{date}         & str       & Fecha ISO-8601 \\
    \texttt{source}       & str       & Fuente (\texttt{xataka\_mobile}, etc.) \\
    \texttt{category}     & str       & Categoría canónica: \texttt{mobile}, \texttt{pc}, \texttt{general} \\
    \texttt{subcategory}  & str       & Etiqueta fina del spider (smartphone, laptop…) \\
    \texttt{brand}        & str       & Marca detectada (Apple, Samsung, Google…) \\
    \texttt{os}           & str       & Sistema operativo (iOS, Android, Windows…) \\
    \texttt{tags}         & list[str] & Etiquetas temáticas \\
    \texttt{metadata}     & dict      & Metadatos adicionales \\
    \bottomrule
  \end{tabular}
\end{table}

% ═══════════════════════════════════════════════════════════════════════════
\section{Modelo de Recuperación Implementado}
% ═══════════════════════════════════════════════════════════════════════════

El modelo de recuperación implementado es \textbf{LSI}
(\textit{Latent Semantic Indexing})~\cite{deerwester1990,manning2008},
que proyecta documentos y consultas en un espacio semántico latente de
dimensión reducida, logrando similitud semántica más allá del solapamiento
exacto de tokens. La arquitectura del sistema se describe en la
sección~\ref{sec:arquitectura}.

\subsection{Arquitectura del sistema}
\label{sec:arquitectura}

\begin{figure}[h]
\centering
\begin{tikzpicture}[
  node distance=0.7cm and 1.8cm,
  box/.style={rectangle, rounded corners=4pt, draw=gray!70,
              fill=blue!6, minimum width=3.2cm, minimum height=0.75cm,
              font=\small, align=center},
  arrow/.style={-{Stealth[length=6pt]}, gray!70, thick},
]

\node[box] (store)   {DocumentStore\\(JSONL)};
\node[box, right=of store] (norm) {TextNormalizer\\tokenize · stem};
\node[box, right=of norm]  (idx)  {InvertedIndex\\posting lists};
\node[box, below=of idx]   (lsi)  {LSIRetriever\\TF-IDF $\to$ SVD};
\node[box, above=of idx]   (bm25) {BM25 search};
\node[box, right=of idx, yshift=-0.75cm] (qp) {QueryProcessor\\expand · filter};
\node[box, right=of qp]    (res)  {Resultados\\rankeados};

\draw[arrow] (store) -- (norm);
\draw[arrow] (norm)  -- (idx);
\draw[arrow] (idx)   -- (bm25);
\draw[arrow] (idx)   -- (lsi);
\draw[arrow] (qp)    -- (res);
\draw[arrow] (bm25)  -| (res);
\draw[arrow] (lsi.east)  -- ++(0.3,0) |- (res.south west);

\end{tikzpicture}
\caption{Arquitectura del sistema SRI. El índice invertido es la estructura
  central compartida por BM25 y LSI.}
\end{figure}

\subsection{Normalización y preprocesado de texto}
\label{sec:normalizer}

La clase \texttt{TextNormalizer} aplica el siguiente pipeline sobre cada
documento y consulta:

\begin{enumerate}
  \item \textbf{Minúsculas}: unificación de casing.
  \item \textbf{Eliminación de URLs}: \texttt{https?://\textbackslash{}S+}.
  \item \textbf{Eliminación de caracteres no alfanuméricos}: se conservan
    letras (incluyendo acentuadas y ñ), dígitos y espacios.
  \item \textbf{Tokenización}: mediante NLTK~\cite{bird2009} \textit{punkt} en modo español.
  \item \textbf{Eliminación de stopwords}: lista combinada español + inglés
    de NLTK~\cite{bird2009} (574 términos).
  \item \textbf{Stemming}: \textit{SnowballStemmer}~\cite{bird2009} en español, que reduce
    las variantes morfológicas a su raíz (``cámara'' $\to$ \textit{cam},
    ``procesadores'' $\to$ \textit{proces}).
\end{enumerate}

El título de cada documento se concatena tres veces antes del contenido para
potenciar su peso en el índice sin modificar las fórmulas de puntuación.

\subsection{Índice Invertido}
\label{sec:index}

\subsubsection{Estructura de datos}

El índice invertido se implementa en la clase \texttt{InvertedIndex} como
estructura central de datos sobre la que opera el recuperador LSI:

\[
  \text{\_index} : \textit{term} \;\longrightarrow\; \{doc\_id : tf\}
\]

donde $tf$ es la frecuencia bruta del término en el documento.  Además se
mantiene por documento:

\[
  \text{\_doc\_info}[d] = \{
    \textit{length},\; \textit{title},\; \textit{url},\;
    \textit{source},\; \textit{date},\; \textit{tags},\; \textit{category}
  \}
\]

Esta estructura de datos es consumida directamente por \texttt{LSIRetriever}
para construir la matriz TF-IDF sin necesidad de re-tokenizar el corpus
(véase sección~\ref{sec:lsi}).

\subsection{Modelo LSI (\textit{Latent Semantic Indexing})}
\label{sec:lsi}

\subsubsection{Motivación}

Los modelos basados en coincidencia exacta de tokens no detectan relaciones
semánticas entre términos. En un corpus con abundantes sinónimos tecnológicos
(``cámara'' / ``camera'' / ``foto'', ``procesador'' / ``chip'' / ``CPU'')
una consulta puede no encontrar documentos semánticamente equivalentes
si no comparten términos exactos. LSI~\cite{deerwester1990,manning2008}
resuelve esto proyectando documentos y consultas en un espacio latente
compartido, donde la similitud refleja relaciones conceptuales en lugar
de solapamiento léxico.

\subsubsection{Construcción de la matriz TF-IDF}

La clase \texttt{LSIRetriever} reutiliza las posting lists del
\texttt{InvertedIndex} para construir la matriz
dispersa $A \in \mathbb{R}^{N \times V}$ con pesos TF-IDF~\cite{manning2008}:

\[
  A_{d,t} = \frac{tf_{t,d}}{|d|} \cdot
    \left(\log\frac{N+1}{df_t+1} + 1\right)
\]

Las filas se normalizan con L2 para que la similaridad coseno sea equivalente
al producto escalar.

\subsubsection{SVD truncada}

Se aplica una descomposición en valores singulares (SVD) truncada~\cite{deerwester1990} de rango $k$, implementada con scikit-learn~\cite{pedregosa2011}:

\[
  A \approx U_k \,\Sigma_k\, V_k^\top
\]

donde $k = \min(\texttt{n\_components},\, N-1,\, V-1)$.  La matriz
$U_k \Sigma_k \in \mathbb{R}^{N \times k}$ contiene los vectores LSI de
los documentos, normalizados con L2.

\subsubsection{Búsqueda en el espacio LSI}

Dada una consulta $q$, su vector TF-IDF se proyecta al espacio LSI~\cite{deerwester1990}
mediante la transformación aprendida por \texttt{TruncatedSVD}~\cite{pedregosa2011}:

\[
  \vec{q}_\text{LSI} = \text{TruncatedSVD.transform}(\vec{q}_\text{tfidf})
\]

La puntuación de cada documento $d_j$ es la similitud coseno:

\[
  \text{score}(d_j, q)
    = \frac{\vec{d}_j \cdot \vec{q}_\text{LSI}}
           {\|\vec{d}_j\|\;\|\vec{q}_\text{LSI}\|}
\]

Al estar ambos vectores normalizados en L2, la operación se reduce a un
producto matricial $D \cdot \vec{q}_\text{LSI}^\top$ vectorizado sobre toda
la colección.

\subsection{Procesador de consultas}
\label{sec:qp}

La clase \texttt{QueryProcessor} aplica el siguiente pipeline antes de
enviar la consulta al recuperador:

\begin{enumerate}
  \item \textbf{Limpieza}: normalización de espacios y minúsculas.
  \item \textbf{Extracción de filtros explícitos}: la sintaxis
    \texttt{clave:valor} (p.\,ej.\ \texttt{category:mobile}) permite
    filtros estructurados incrustados en la consulta de texto libre.
  \item \textbf{Inferencia de categoría}: si no hay filtro explícito, se
    detecta la categoría (\texttt{mobile} / \texttt{pc}) a partir de un
    diccionario de indicadores léxicos (smartphones, laptops, etc.).
  \item \textbf{Expansión con sinónimos}: se añaden sinónimos del dominio
    tecnológico para mejorar el recall.
    Ejemplo:
    \begin{center}
      \texttt{"iphone camara nocturna"}
      $\;\longrightarrow\;$
      \texttt{"iphone camara nocturna camera foto fotografia"}
    \end{center}
\end{enumerate}

Opcionalmente, los resultados pueden ser reordenados por fecha
(\texttt{rerank\_by\_date}) o por diversidad mediante
\textbf{MMR} (\textit{Maximal Marginal Relevance})~\cite{carbonell1998}:

\[
  \text{MMR}(d_i) = (1-\lambda)\,\text{Rel}(d_i)
                  - \lambda \max_{d_j \in S}\,\text{Sim}(d_i, d_j)
\]

donde $S$ es el conjunto de documentos ya seleccionados y
$\text{Sim}(\cdot,\cdot)$ es la similaridad de Jaccard sobre etiquetas.

\subsection{Pipeline completo de recuperación}

\begin{figure}[h]
\centering
\begin{tikzpicture}[
  node distance=0.5cm,
  box/.style={rectangle, rounded corners=4pt, draw=gray!60,
              fill=blue!5, minimum width=4.8cm, minimum height=0.7cm,
              font=\small, align=center},
  arrow/.style={-{Stealth[length=6pt]}, gray!60, thick},
]
\node[box] (q0) {Consulta en lenguaje natural};
\node[box, below=of q0] (qp) {QueryProcessor\\normaliza · expande · infiere categoría};
\node[box, below=of qp] (lsi) {LSIRetriever.retrieve(pq.text,\\category\_filter=pq.category)};
\node[box, below=of lsi] (rank) {Ranking por similitud coseno LSI};
\node[box, below=of rank] (out) {Top-$k$ resultados};

\draw[arrow] (q0) -- (qp);
\draw[arrow] (qp)  -- (lsi);
\draw[arrow] (lsi) -- (rank);
\draw[arrow] (rank)-- (out);
\end{tikzpicture}
\caption{Pipeline completo de una consulta en el sistema SRI.}
\end{figure}

\subsection{Persistencia de índices}
\label{sec:persist}

El modelo LSI se serializa a disco para evitar reentrenamientos:

\begin{table}[h]
  \centering
  \caption{Archivos de persistencia del modelo LSI.}
  \begin{tabular}{lll}
    \toprule
    \textbf{Archivo} & \textbf{Formato} & \textbf{Contenido} \\
    \midrule
    \texttt{indexes/index/inverted\_index.pkl} & pickle & posting lists, \texttt{doc\_info} \\
    \texttt{indexes/index/index\_meta.json}    & JSON   & vocab, $N$ \\
    \texttt{indexes/lsi/lsi\_model.pkl}        & pickle & SVD, matriz doc, vocab, IDF \\
    \texttt{indexes/lsi/lsi\_meta.json}        & JSON   & varianza explicada, $k$, $N$ \\
    \bottomrule
  \end{tabular}
\end{table}

% ═══════════════════════════════════════════════════════════════════════════
\section{Extracción de datos}
% ═══════════════════════════════════════════════════════════════════════════

\subsection{Spiders implementados}

Los spiders se implementan con Scrapy~\cite{scrapy2023} heredando de la clase
base abstracta \texttt{Extract}, que garantiza un esquema unificado y buenas
prácticas de \textit{crawling} (respeto de \texttt{robots.txt},
\textit{AutoThrottle}, caché HTTP).

\begin{table}[h]
  \centering
  \caption{Spiders de extracción implementados.}
  \begin{tabular}{lll}
    \toprule
    \textbf{Spider} & \textbf{Fuente} & \textbf{Ítem} \\
    \midrule
    \texttt{xataka\_mobile}  & xataka.com (móviles) & \texttt{MobileItem} \\
    \texttt{xataka\_pc}      & xataka.com (PC)      & \texttt{PCItem} \\
    \bottomrule
  \end{tabular}
\end{table}

\subsection{Pipelines de Scrapy}

Los ítems extraídos pasan por tres pipelines en orden:

\begin{enumerate}
  \item \textbf{TimestampPipeline} (prioridad 100): añade el campo
    \texttt{scraped\_at} con marca de tiempo UTC.
  \item \textbf{DuplicatesPipeline} (prioridad 300): descarta ítems cuya
    URL ya fue vista en la sesión actual.
  \item \textbf{JsonStoragePipeline} (prioridad 500): persiste los ítems
    como JSONL en \texttt{data/\{mobile,pc,general\}/spider\_YYYY-MM-DD.jsonl}.
\end{enumerate}

% ═══════════════════════════════════════════════════════════════════════════
% References (LNCS style: author surname, initials: title. venue, pages (year))
% ═══════════════════════════════════════════════════════════════════════════

\begin{thebibliography}{6}

\bibitem{deerwester1990}
Deerwester, S., Dumais, S.T., Furnas, G.W., Landauer, T.K., Harshman, R.:
Indexing by Latent Semantic Analysis.
J. Am. Soc. Inf. Sci. \textbf{41}(6), 391--407 (1990).
\url{https://doi.org/10.1002/(SICI)1097-4571(199009)41:6<391::AID-ASI1>3.0.CO;2-9}

\bibitem{manning2008}
Manning, C.D., Raghavan, P., Schütze, H.:
Introduction to Information Retrieval.
Cambridge University Press, Cambridge (2008).
\url{https://nlp.stanford.edu/IR-book/}

\bibitem{carbonell1998}
Carbonell, J., Goldstein, J.:
The use of MMR, diversity-based reranking for reordering documents and
producing summaries.
In: Proc. 21st ACM SIGIR, pp. 335--336 (1998).
\url{https://doi.org/10.1145/290941.291025}

\bibitem{bird2009}
Bird, S., Klein, E., Loper, E.:
Natural Language Processing with Python.
O'Reilly Media, Sebastopol (2009).
\url{https://www.nltk.org/book/}

\bibitem{pedregosa2011}
Pedregosa, F., et al.:
Scikit-learn: Machine Learning in Python.
J. Mach. Learn. Res. \textbf{12}, 2825--2830 (2011).
\url{https://scikit-learn.org/}

\bibitem{scrapy2023}
Scrapy developers:
Scrapy documentation (2023).
\url{https://docs.scrapy.org/}

\end{thebibliography}

\end{document}
